% Part: sets-functions-relations
% Chapter: induction
% Section: induction

\documentclass[../../../include/open-logic-section]{subfiles}

\begin{document}

\olfileid{sfr}{ind}{int}

\olsection{Pambuka}

\begin{explain}
Induksi minangka teknik bukti kang kerep digunakake ing matematika lan
logika. Mbokmenawa sampeyan wis ngerti saka matematika, nalika
nggunakake induksi marang wilangan asli. Pranyata akeh jinis
himpunan objek matematika lan logika kang bisa ditangani nganggo induksi.

Umume, induksi ngidini kita mbuktekake pratelan universal, yaiku
nuduhake yen saben objek saka sawijining jinis nduweni sipat tartamtu.
Mligine, induksi asring migunani nalika cara bukti ``introduksi universal''
kakehan kerumitane. Induksi mung lumaku kanggo objek matematika kang
dibangun kanthi cara khusus: yen saben unsur ing himpunan minangka unsur
dhasar utawa dibangun saka unsur dhasar, mula kita bisa nggunakake
induksi marang himpunan kasebut. Kanthi luwih tliti:
\end{explain}

\begin{defn}
Himpunan $S$ \emph{katutup} tumrap fungsi $f$ yen lan mung yen
saben $x \in S$ ndadekake $f(x) \in S$.
Kanggo fungsi kanthi $n \ge 1$ masukan, katutup ateges
$f(x_1,\ldots,x_n) \in S$ saben $x_1,\ldots,x_n \in S$.
\end{defn}

\begin{explain}
Contone, wilangan asli ($\mathbb{N}$) katutup tumrap fungsi penerus
$f(x) = x+1$. Yen $x$ minangka wilangan asli, semono uga $x+1$.
Nanging wilangan asli ora katutup tumrap fungsi sadurunge
$f(x) = x-1$, amarga $0$ minangka wilangan asli, dene $-1$ dudu.
\end{explain}

\begin{defn}
Sipat $P$ diarani \emph{dijaga tumrap} fungsi $f(x)$ yen kanggo
saben objek $a$ ing domain $f$, $P(a)$ ndadekake $P(f(a))$
(yen $a$ nduweni sipat $P$, mula $f(a)$ uga nduweni sipat kasebut).
Kanggo fungsi kanthi $n \ge 1$ masukan, dijaga ateges
$P(a_1),\ldots,P(a_n)$ ndadekake $P(f(a_1,\ldots,a_n))$.
\end{defn}

\begin{explain}
Sipat ``minangka wilangan genep'' dijaga tumrap fungsi $f(x) = x+2$,
amarga yen $x$ genep, mula $x+2$ uga genep. Sipat minangka wilangan
genep ora dijaga tumrap fungsi penerus, amarga yen $x$ genep,
$x+1$ ganjil.

Induksi lumaku kanggo wilangan asli amarga saben wilangan asli bisa
digayuh saka 0 kanthi ngetrapake fungsi penerus kaping winates.
Iki minangka ciri saben himpunan kang bisa ditangani nganggo induksi.
\end{explain}

\begin{defn}[Induksi marang $\Nat$]
Yen 0 nduweni sipat $P$, lan yen $P$ dijaga tumrap fungsi penerus,
mula saben wilangan asli nduweni sipat $P$.
\end{defn}

\begin{ex} Gunggung wilangan bulat saka 1 nganti $n$ yaiku $\frac{n(n+1)}{2}$. \\

\textbf{Kasus dhasar.} Gunggung kosong nalika n padha karo 0 yaiku
$0=\frac{0\cdot 1}{2}$.\\

\textbf{Langkah induktif.} Upamane sipat kasebut lumaku kanggo $k$.
Kita bakal nuduhake yen lumaku kanggo $k+1$. Tegese, upamane $P(k)$
lumaku, yaiku,
\[ 0 + 1 + \ldots + k = \frac{k(k+1)}{2} \]
Kita bakal nuduhake yen $P(k+1)$ lumaku, yaiku
\[ 0 + 1 + \ldots + k + (k+1) = \frac{(k+1)(k+2)}{2} \]
kanthi nggunakake anggepan $P(k)$.

\begin{align*}
0 + 1 + \ldots + k &= \frac{k(k+1)}{2} \tag{Anggepan} \\
0 + 1 + \ldots + k + (k+1) &= \frac{k(k+1)}{2} + (k+1)
\tag{Tambah $k+1$ ing loro sisih} \\
&= \frac{k(k+1) + 2(k+1)}{2} \\
&= \frac{k^2 + k + 2k +2}{2} \\
&=\frac{(k+1)(k+2)}{2}
\end{align*}
Dadi langkah induktif kasil, lan kita wis mbuktekake pratelan kasebut.
\end{ex}

\begin{explain}
Kanggo !!{formula}s, unsur dhasare yaiku !!{formula}s atomik.
!!{formula}s kang luwih rumit diolehake kanthi nggabungake
!!{formula}s kang luwih prasaja nganggo operator. Saben operator
bisa dianggep cocog karo fungsi kanthi cacah !!{formula}s kang
jumbuh karo aritase, kang ngasilake !!{formula} anyar kanthi
ngetrapake operator kasebut. Contone, fungsi kang nggabungake
loro !!{formula}s $!A$ lan $!B$ dadi konjungsi bisa ditulis
minangka $\mathcal E _\land (!A, !B) = !A \land !B$.
Kita bisa ngarani fungsi kasebut \emph{fungsi pambangun formula}.
Himpunan !!{formula}s katutup tumrap fungsi kasebut:
saben $!A$ lan $!B$ minangka !!{formula}s, mula $\lnot !A$,
$!A \land !B$, lan sapanunggalane uga minangka formula.
\end{explain}

\begin{defn}[Induksi marang !!{formula}s]
Yen saben !!{formula} atomik nduweni sipat $P$ lan $P$ dijaga
tumrap fungsi pambangun !!{formula}, mula saben !!{formula}
nduweni sipat $P$.
\end{defn}

\begin{explain}
Definisi iki menehi ``cara'' kanggo bukti induktif marang
!!{formula}s. Yen kita dijaluk nuduhake yen saben !!{formula}
nduweni sipat $P$, tindakna langkah ing ngisor iki:\\

\textbf{Kasus dhasar.} Upamane $!A$ minangka !!{formula} atomik.
[\ldots] Mula $!A$ nduweni sipat $P$.

\textbf{Langkah induktif.} Upamane $!A$ lan $!B$ minangka
!!{formula}s kang kalorone nduweni sipat $P$.

Kasus $\lnot$: [\ldots] Mula $\lnot !A$ nduweni sipat $P$.

Kasus $\land$: [\ldots] Mula $!A \land !B$ nduweni sipat $P$.

Kasus $\lor$: [\ldots] Mula $!A \lor !B$ nduweni sipat $P$.

Kasus $\lif$: [\ldots] Mula $!A \lif !B$ nduweni sipat $P$.

Kasus $\liff$: [\ldots] Mula $!A \liff !B$ nduweni sipat $P$.

Kasus $\lforall[]$: [\ldots] Mula $\lforall[x] !A$ nduweni sipat $P$.

Kasus $\lexists[]$: [\ldots] Mula $\lexists[x] !A$ nduweni sipat $P$. \\

Mula saben !!{formula} nduweni sipat $P$.
\end{explain}

\end{document}
